\documentclass[spanish,letterpaper,oneside]{article}

\def\documenttitle {El lugar perfecto · Ficha de trabajo}
\def\documentsubtitle {Temas selectos de matemáticas I · Cohorte SD26}
\def\documentsubject {Actividad 2 · Material para estudiantes}
\def\documentauthor {Mtro. Martín Jonathan de la Cruz Muñoz}
\def\coursename {Temas selectos de matemáticas I}
\def\coursecode {PFM02}
\def\universityname {Instituto de Investigación, Innovación y Estudios de Posgrado para la Educación del Estado de Nuevo León}
\def\universityfaculty {Maestría en Aprendizaje y Enseñanza de las Matemáticas}
\def\universitydepartment {Temas selectos de matemáticas I}
\def\universitydepartmentimage {img/iiiepe/marca-iiiepe.png}
\def\universitydepartmentimagecfg {height=1.75cm}
\def\universitylocation {Monterrey, Nuevo León}
\def\studentfullname {Mtro. Martín Jonathan de la Cruz Muñoz}
\def\studentid {00874}
\def\studentemail {martin.delacruz@campus.iiiepe.edu.mx}
\def\programperiod {2do periodo, agosto--noviembre de 2026 (SD26)}
\def\programmodality {Modalidad presencial y en línea}
\def\programshort {MAEM}

\def\authortable {
  \begin{tabular}{ll}
    \textbf{Docente:} & \studentfullname \\
    Matrícula: & \studentid \\
    Unidad: & \coursename \\
    Periodo: & \programperiod \\
    \multicolumn{2}{l}{Fecha: 17 de octubre de 2026}
  \end{tabular}
}

\input{template}

\usepackage{helvet}
\renewcommand{\familydefault}{\sfdefault}
\usepackage{amsmath,amssymb,array,longtable,tikz,tabularx,colortbl,enumitem}
\definecolor{iiiepeNavy}{HTML}{163B5C}
\definecolor{iiiepeBlue}{HTML}{1F6F9F}
\definecolor{iiiepePale}{HTML}{EEF5F8}
\def\portraittitlecolor {iiiepeNavy}

\setlength{\parskip}{5pt}\setlength{\parindent}{0pt}
\begin{document}\templatePortrait\templatePagecfg\templateFinalcfg\fancyfoot{}\fancyfoot[L]{\scriptsize IIIEPE · MAEM · SD26}\fancyfoot[C]{\scriptsize\thepage}\fancyfoot[R]{\scriptsize Ficha Act. 2}\setstretch{1.2}\raggedbottom

\section*{Construimos y comparamos lugares geométricos}
Nombre: \rule{8cm}{.2pt}\hfill Fecha: \rule{2.5cm}{.2pt}

\subsection*{Actividad 1. Regla, compás y mediatrices}
Dado un triángulo $ABC$, traza con regla y compás las mediatrices de dos lados para ubicar el punto $O$ y verifica con tu compás la circunferencia que pasa por los tres vértices. ¿Por qué basta trazar dos mediatrices?
\vspace{2.0cm}

\subsection*{Actividad 2. Registro de exploración en GeoGebra}
Completa la tabla tras arrastrar los vértices en GeoGebra Classic:

\begin{center}
\begin{tabular}{|l|p{5.2cm}|p{5.2cm}|}\hline
\textbf{Tipo de triángulo} & \textbf{Ubicación de $O$ (Circuncentro)} & \textbf{Ubicación de $I$ (Incentro)} \\\hline
Acutángulo & & \\[0.6cm]\hline
Rectángulo & & \\[0.6cm]\hline
Obtusángulo & & \\[0.6cm]\hline
\end{tabular}
\end{center}

\end{document}
